% AI tự sinh cho Hội đồng AI · Lớp 10

\begin{ex}%[0D0H1-1]
% nguon: Hội đồng AI (Gemini soạn · DeepSeek giải · DeepSeek phản biện · Gemini duyệt)
Một hộp chứa $5$ quả cầu màu xanh và $4$ quả cầu màu đỏ có kích thước và khối lượng như nhau. Lấy ngẫu nhiên đồng thời $3$ quả cầu từ hộp. Xét các biến cố sau:
\begin{itemize}
    \item $A$: "Trong ba quả cầu lấy ra, có ít nhất một quả cầu màu đỏ";
    \item $B$: "Cả ba quả cầu lấy ra đều có màu xanh";
    \item $C$: "Trong ba quả cầu lấy ra, có nhiều nhất hai quả cầu màu xanh".
\end{itemize}
Khẳng định nào sau đây là \textbf{sai}?
\choice
{$A$ và $B$ là hai biến cố đối nhau}
{$A$ và $C$ là hai biến cố bằng nhau}
{$B$ và $C$ là hai biến cố đối nhau}
{\True $A$ và $C$ là hai biến cố xung khắc}
\loigiai{
Tổng số quả cầu lấy ra là $3$.
\begin{itemize}
    \item Biến cố $A$ xảy ra khi số quả cầu màu đỏ lấy ra nhận một trong các giá trị $\{1; 2; 3\}$, tương ứng số quả cầu màu xanh nhận một trong các giá trị $\{2; 1; 0\}$, nghĩa là có nhiều nhất $2$ quả cầu màu xanh. Do đó, $A$ và $C$ là hai biến cố bằng nhau ($A = C$).
    \item Biến cố đối của $A$ là $\overline{A}$: "Không có quả cầu màu đỏ nào được lấy ra", tức là "Cả ba quả cầu lấy ra đều có màu xanh", chính là biến cố $B$. Do đó $\overline{A} = B$, suy ra $A$ và $B$ là hai biến cố đối nhau. Vì $A = C$ nên $B$ và $C$ cũng là hai biến cố đối nhau.
    \item Do hộp có $4$ quả cầu đỏ và $5$ quả cầu xanh nên ta hoàn toàn có thể lấy ra $3$ quả cầu thoả mãn có ít nhất một quả cầu đỏ (chẳng hạn $1$ quả đỏ và $2$ quả xanh), tức là $A \neq \emptyset$. Vì $A = C$ và $A \neq \emptyset$ nên $A \cap C = A \neq \emptyset$. Do đó, $A$ và $C$ không phải là hai biến cố xung khắc.
\end{itemize}
Vậy khẳng định "$A$ và $C$ là hai biến cố xung khắc" là khẳng định sai.
}
\end{ex}

\begin{ex}%[0D0H1-1]
% nguon: Hội đồng AI (Gemini soạn · DeepSeek giải · DeepSeek phản biện · Gemini duyệt)
Gieo đồng thời hai con xúc xắc cân đối và đồng chất. Gọi $\Omega$ là không gian mẫu của phép thử. Xét hai biến cố:
\begin{itemize}
    \item $A$: "Tổng số chấm xuất hiện trên hai con xúc xắc chia hết cho $4$";
    \item $B$: "Có ít nhất một con xúc xắc xuất hiện mặt $2$ chấm".
\end{itemize}
Số phần tử của biến cố giao $A \cap B$ là
\choice
{$2$}
{\True $3$}
{$4$}
{$5$}
\loigiai{
Biến cố giao $A \cap B$ là biến cố: "Tổng số chấm xuất hiện chia hết cho $4$ và có ít nhất một con xúc xắc xuất hiện mặt $2$ chấm".

Gọi kết quả của phép thử là cặp số $(x; y)$ với $x, y \in \{1; 2; 3; 4; 5; 6\}$.

Tổng số chấm $S = x + y \in \{2; 3; \ldots; 12\}$. Do $S$ chia hết cho $4$ nên $S \in \{4; 8; 12\}$.
\begin{itemize}
    \item Trường hợp $S = 4$: Kết hợp điều kiện có ít nhất một số bằng $2$, ta có cặp kết quả $(2; 2)$.
    \item Trường hợp $S = 8$: Kết hợp điều kiện có ít nhất một số bằng $2$, ta có các cặp kết quả $(2; 6)$ và $(6; 2)$.
    \item Trường hợp $S = 12$: Chỉ có cặp $(6; 6)$, không có con xúc xắc nào xuất hiện mặt $2$ chấm.
\end{itemize}
Vậy biến cố $A \cap B = \{(2; 2); (2; 6); (6; 2)\}$, suy ra số phần tử của biến cố $A \cap B$ là $n(A \cap B) = 3$.
}
\end{ex}

\begin{ex}%[0D0H1-1]
% nguon: Hội đồng AI (Gemini soạn · DeepSeek giải · DeepSeek phản biện · Gemini duyệt)
Gieo một đồng xu cân đối và đồng chất ba lần liên tiếp. Kí hiệu $S$ là kết quả xuất hiện mặt sấp, $N$ là kết quả xuất hiện mặt ngửa. Xét hai biến cố:
\begin{itemize}
    \item $A$: "Có đúng một lần xuất hiện mặt ngửa";
    \item $B$: "Lần gieo thứ nhất xuất hiện mặt sấp".
\end{itemize}
Cho các khẳng định sau:
\begin{enumerate}[(I)]
    \item Không gian mẫu $\Omega$ có $8$ phần tử.
    \item Tập hợp mô tả biến cố $A$ là $A = \{SSN; SNS; NSS\}$.
    \item Biến cố giao $A \cap B = \{SSN; SNS\}$.
    \item Hai biến cố $A$ và $B$ xung khắc với nhau.
\end{enumerate}
Trong các khẳng định trên, có bao nhiêu khẳng định đúng?
\choice
{$1$}
{$2$}
{\True $3$}
{$4$}
\loigiai{
Mỗi lần gieo có $2$ kết quả có thể xảy ra ($S$ hoặc $N$). Khi gieo $3$ lần liên tiếp, số phần tử của không gian mẫu là $n(\Omega) = 2^3 = 8$. Do đó khẳng định (I) đúng.

Biến cố $A$ là "Có đúng một lần xuất hiện mặt ngửa", các kết quả thuận lợi cho $A$ là:
$$A = \{SSN; SNS; NSS\}.$$
Do đó khẳng định (II) đúng.

Biến cố $B$ là "Lần gieo thứ nhất xuất hiện mặt sấp", tức là:
$$B = \{SSS; SSN; SNS; SNN\}.$$
Khi đó, biến cố giao $A \cap B$ gồm các kết quả vừa thuộc $A$ vừa thuộc $B$:
$$A \cap B = \{SSN; SNS\}.$$
Do đó khẳng định (III) đúng.

Vì $A \cap B = \{SSN; SNS\} \neq \emptyset$ nên $A$ và $B$ không phải là hai biến cố xung khắc. Do đó khẳng định (IV) sai.

Vậy có tất cả $3$ khẳng định đúng.
}
\end{ex}
